\documentclass[11pt]{article}
\usepackage[margin=1in]{geometry}
\usepackage{amsmath,amssymb,amsthm}
\usepackage{hyperref}

\newtheorem{theorem}{Theorem}
\newtheorem{lemma}[theorem]{Lemma}
\newtheorem{proposition}[theorem]{Proposition}
\newtheorem{corollary}[theorem]{Corollary}
\theoremstyle{definition}
\newtheorem{definition}[theorem]{Definition}
\theoremstyle{remark}
\newtheorem{remark}[theorem]{Remark}

\newcommand{\Q}{\mathbb Q}
\newcommand{\R}{\mathbb R}
\newcommand{\Z}{\mathbb Z}
\newcommand{\N}{\mathbb N}
\newcommand{\id}{\mathrm{id}}
\newcommand{\Fix}{\operatorname{Fix}}

\title{Disproof of Conjecture 00000007608:\\
a rotation has infinitely many decompositions into two reflections}
\author{jilint777}
\date{2026-10-04}

\begin{document}
\sloppy
\maketitle

\begin{abstract}
Conjecture 00000007608 asserts that the composition of two reflections is a rotation, that the
decomposition (of a rotation into two reflections) is unique, and that the exception to
uniqueness is the case of parallel reflections. The uniqueness clause is false for
\emph{every} rotation $\rho\neq\id$ of the plane. Writing $R_\alpha$ for the reflection in the
line through the centre at angle $\alpha$, we have $\rho_\theta=R_\beta\circ R_\alpha$ whenever
$\beta-\alpha\equiv\theta/2\pmod\pi$, so $\alpha$ can be chosen freely. None of these
decompositions has parallel mirrors. The simplest instance is
$\left(\begin{smallmatrix}0&1\\1&0\end{smallmatrix}\right)
\left(\begin{smallmatrix}1&0\\0&-1\end{smallmatrix}\right)=
\left(\begin{smallmatrix}0&-1\\-1&0\end{smallmatrix}\right)
\left(\begin{smallmatrix}-1&0\\0&1\end{smallmatrix}\right)=
\left(\begin{smallmatrix}0&-1\\1&0\end{smallmatrix}\right)$. Its four mirrors are the $x$-axis,
the line $y=x$, the $y$-axis and the line $y=-x$. In Lean~4 (core library only) we formalise
rational $2\times2$ reflections and rotations from scratch. We prove that the decompositions of
any rational rotation $\rho$ are exactly the pairs $(A,\rho A)$. We prove that for every rational
$\rho\neq\id$ the Pythagorean family $A_t$ ($t\in\N$) gives infinitely many non-parallel
decompositions with pairwise different mirrors, no three of which share an unordered mirror
pair. The uniqueness clause is refuted in its weakest literal form (unordered pairs of mirrors).
We also treat $\R^3$ and affine reflections, and we discuss the one charitable reading under
which the clause becomes true (``unique up to rotating both mirrors'').
\end{abstract}

\section{The claim}

The English statement reads: \emph{``Definition: Reflections as hyperplane mirror structure.
Conjecture: The composition of two reflections is a rotation and the decomposition is unique;
the exception of uniqueness is parallel reflections; and the exception of parallel reflections
is translation composition.''} The Chinese version says the same thing: the composition of two
reflections is a rotation and the decomposition is unique (``fen jie wei yi''); the exception
to uniqueness is parallel reflections; the exception for parallel reflections is a composite
translation. We split it into four clauses:
\begin{itemize}
\item[(C1)] the composition of two reflections is a rotation;
\item[(C2)] the decomposition of a rotation into two reflections is unique;
\item[(C3)] the exception to (C2) is the case of parallel mirrors;
\item[(C4)] the exception to (C1) is the case of parallel mirrors, whose composition is a
translation.
\end{itemize}
Clauses (C2) and (C3) together form the \emph{uniqueness clause}:
\begin{quote}
(U) If $\rho$ is a rotation and $\rho=B\circ A=B'\circ A'$ with reflections $A,B,A',B'$, and
neither pair has parallel mirrors, then the two decompositions coincide.
\end{quote}
We refute (U). Because the conjecture is a conjunction, this refutes the conjecture. Clauses (C1)
and (C4) are discussed in Section~\ref{sec:readings}; the disproof does not depend on them.

\section{Set-up}

A \emph{reflection} of $\R^2$ (fixing the origin) is an orthogonal matrix $A$ with $\det A=-1$
and $A^2=I$. A \emph{rotation} is an orthogonal matrix with $\det=+1$. The \emph{mirror} of a
reflection $A$ is its fixed line $\Fix A=\ker(A-I)$. Two mirrors through the origin are parallel
iff they are equal. For $\alpha,\theta\in\R$ put
\[
R_\alpha=\begin{pmatrix}\cos2\alpha&\sin2\alpha\\ \sin2\alpha&-\cos2\alpha\end{pmatrix},\qquad
\rho_\theta=\begin{pmatrix}\cos\theta&-\sin\theta\\ \sin\theta&\cos\theta\end{pmatrix}.
\]
$R_\alpha$ is the reflection in the line $\R(\cos\alpha,\sin\alpha)$, and $R_\alpha=R_{\alpha'}$
iff $\alpha\equiv\alpha'\pmod\pi$.

\begin{lemma}\label{lem:normal}
Over any ordered field $K$ (in particular $\Q$ or $\R$), an orthogonal $2\times2$ matrix with
determinant $-1$ has the form $\left(\begin{smallmatrix}a&b\\ b&-a\end{smallmatrix}\right)$ with
$a^2+b^2=1$, and it is automatically an involution. An orthogonal matrix with determinant $+1$
has the form $\left(\begin{smallmatrix}x&-y\\ y&x\end{smallmatrix}\right)$ with $x^2+y^2=1$. Over
$\R$ these are exactly the $R_\alpha$ and the $\rho_\theta$.
\end{lemma}
\begin{proof}
Let $M=\left(\begin{smallmatrix}a&b\\ c&d\end{smallmatrix}\right)$ with $MM^{T}=I$, so
$a^2+b^2=c^2+d^2=1$. If $ad-bc=-1$, then
$(b-c)^2+(a+d)^2=(a^2+b^2)+(c^2+d^2)+2(ad-bc)=0$, so $c=b$ and $d=-a$. Then
$M^2=(a^2+b^2)I=I$. If $ad-bc=1$, then $(b+c)^2+(a-d)^2=2-2=0$ gives $c=-b$ and $d=a$.
\end{proof}

\begin{lemma}\label{lem:product}
$R_\beta R_\alpha=\rho_{2(\beta-\alpha)}$ and $\rho_\theta R_\alpha=R_{\alpha+\theta/2}$.
\end{lemma}
\begin{proof}
By direct multiplication and the addition formulas,
\[
R_\beta R_\alpha=\begin{pmatrix}\cos2\beta\cos2\alpha+\sin2\beta\sin2\alpha&
\cos2\beta\sin2\alpha-\sin2\beta\cos2\alpha\\ \sin2\beta\cos2\alpha-\cos2\beta\sin2\alpha&
\sin2\beta\sin2\alpha+\cos2\beta\cos2\alpha\end{pmatrix}=\rho_{2\beta-2\alpha}.
\]
The second identity follows by multiplying the first, with $\beta=\alpha+\theta/2$, on the right
by $R_\alpha=R_\alpha^{-1}$.
\end{proof}

\section{The main theorem}

\begin{theorem}\label{thm:main}
Let $\rho=\rho_\theta\neq I$ be a rotation of $\R^2$.
\begin{enumerate}
\item The decompositions $\rho=B\circ A$ into two reflections are exactly the pairs
$(A,\rho A)$ with $A$ an arbitrary reflection, i.e.\ $(R_\alpha,R_{\alpha+\theta/2})$ for
$\alpha\in\R/\pi\Z$.
\item In none of them are the two mirrors parallel.
\item Distinct $\alpha\in\R/\pi\Z$ give distinct ordered pairs and distinct first mirrors. Each
unordered pair of mirrors arises for at most two values of $\alpha$. Hence there are
uncountably many decompositions, even counted as unordered pairs of mirrors.
\end{enumerate}
\end{theorem}
\begin{proof}
(1) If $BA=\rho$, then $B=BAA=\rho A$. Conversely, if $A=R_\alpha$, then
$\rho A=R_{\alpha+\theta/2}$ is a reflection (Lemma~\ref{lem:product}), and $(\rho A)A=\rho A^2=\rho$.
(2) If $A$ and $B=\rho A$ had the same mirror, a nonzero vector $v$ on it would satisfy
$Av=v=Bv=\rho Av=\rho v$. A rotation $\neq I$ fixes no nonzero vector, since
$\det(\rho_\theta-I)=2-2\cos\theta\neq0$. Equivalently, the mirrors are at angle
$\theta/2\not\equiv0\pmod\pi$.
(3) The first mirror determines $\alpha$ modulo $\pi$. An unordered pair $\{m,m'\}$ contains at
most two lines, so at most two decompositions have their first mirror in it.
\end{proof}

So uniqueness fails for \emph{every} nontrivial rotation. The exception named in (C3), parallel
mirrors, never occurs for $\rho\neq I$: parallel (equal) linear mirrors give $A=B$ and
$BA=I$. The only rotation that has decompositions with parallel mirrors is the identity, and
for it \emph{all} decompositions $(A,A)$ have equal mirrors.

\section{The rational version (formalised in Lean)}

Lean's core library has no real numbers, so the formal proof works over $\Q$. That is enough,
because the counterexamples are rational. A rational matrix is stored as $m/q$, with $m$ an
integer matrix and $q>0$, and all predicates are the rational ones with the denominators
cleared. Lemma~\ref{lem:normal} holds verbatim over $\Q$, where the last step uses that a sum
of two rational squares vanishes only if both do. It is formalised as \texttt{isRefl\_iff},
\texttt{rot\_form}, \texttt{refl\_of\_form} and \texttt{rot\_of\_form}.

\begin{proposition}[general statements, all rational rotations]\label{prop:general}
Let $\rho$ be a rational rotation and $A,B$ rational reflections.
\begin{enumerate}
\item $BA$ is a rotation (\texttt{refl\_mul\_refl}; this is (C1) in the linear case).
\item $\rho A$ is a reflection and $(\rho A)A=\rho$ (\texttt{rot\_mul\_refl},
\texttt{decomp\_of\_refl}).
\item If $BA=\rho$, then $B=\rho A$ (\texttt{decomp\_classification}).
\end{enumerate}
\end{proposition}
\begin{proof}
Write $\rho=\left(\begin{smallmatrix}x&-y\\ y&x\end{smallmatrix}\right)$ with $x^2+y^2=1$ and
$A=\left(\begin{smallmatrix}a&b\\ b&-a\end{smallmatrix}\right)$ with $a^2+b^2=1$. Then
$\rho A=\left(\begin{smallmatrix}a'&b'\\ b'&-a'\end{smallmatrix}\right)$ with $a'=xa-yb$ and
$b'=xb+ya$. Moreover $a'^2+b'^2=(x^2+y^2)(a^2+b^2)=1$, so $\rho A$ is a reflection by
Lemma~\ref{lem:normal}. Also $(\rho A)A=\rho(a^2+b^2)=\rho$. Parts (1) and (3) are similar
polynomial identities, given in the Lean file with the denominators cleared.
\end{proof}

\paragraph{The Pythagorean family.} For $t\in\N$ let
\[
A_t=\frac1{1+t^2}\begin{pmatrix}1-t^2&2t\\ 2t&t^2-1\end{pmatrix}.
\]
Since $(1-t^2)^2+(2t)^2=(1+t^2)^2$, Lemma~\ref{lem:normal} shows that $A_t$ is a reflection
(\texttt{Aref\_refl}). Moreover $A_t(1,t)^{T}=(1,t)^{T}$ (\texttt{Aref\_fix}), so the mirror of
$A_t$ is the line of slope $t$. If $A_s$ fixes $(1,t)$, then the second coordinate gives
$2s+(s^2-1)t=(1+s^2)t$, i.e.\ $s=t$ (\texttt{Aref\_fix\_eq}). Hence the mirrors of the $A_t$ are
pairwise different (\texttt{Aref\_distinct}), and so are the matrices (\texttt{Aref\_injective}).
Examples: $A_0=\operatorname{diag}(1,-1)$, $A_1=\left(\begin{smallmatrix}0&1\\
1&0\end{smallmatrix}\right)$, $A_2=\frac15\left(\begin{smallmatrix}-3&4\\
4&3\end{smallmatrix}\right)$.

\begin{lemma}[\texttt{nonparallel}]\label{lem:nonpar}
If $\rho$ is a rational rotation with $\rho\neq I$, then $A_t$ and $\rho A_t$ have different
mirrors.
\end{lemma}
\begin{proof}
Otherwise $\rho A_t$ fixes $v=(1,t)$. Since $A_tv=v$, this gives $\rho v=v$, that is,
$x-yt=1$ and $y+xt=t$. Multiplying the first equation by $t$ and subtracting the second gives
$y(1+t^2)=0$, so $y=0$ and $x=1$, i.e.\ $\rho=I$.
\end{proof}

\begin{theorem}[\texttt{infinitely\_many\_decompositions}, \texttt{not\_uniqueFor}]
\label{thm:rational}
Let $\rho\neq I$ be a rational rotation, and let $D_t=(A_t,\rho A_t)$ for $t\in\N$. Then:
\begin{enumerate}
\item every $D_t$ is a decomposition $\rho=(\rho A_t)\circ A_t$ with non-parallel mirrors;
\item for $s\neq t$ the first mirrors of $D_s$ and $D_t$ differ;
\item for pairwise distinct $s,t,u$, the unordered mirror pairs of $D_s$, $D_t$, $D_u$ are not
all equal.
\end{enumerate}
In particular (U) fails for $\rho$, even when decompositions are compared as unordered pairs of
mirrors.
\end{theorem}
\begin{proof}
(1) is Proposition~\ref{prop:general}(2) together with Lemma~\ref{lem:nonpar}. (2) is
\texttt{Aref\_distinct}. For (3), suppose all three unordered pairs coincide. Write $m_t$ for the
mirror of $A_t$ and $m'_t$ for that of $\rho A_t$. Since $m_t\neq m_s$, equality of
$\{m_s,m'_s\}$ and $\{m_t,m'_t\}$ forces $m_t=m'_s$; likewise $m_u=m'_s$. Hence $m_t=m_u$,
contradicting (2). Applying (U) to $D_0,D_1$ and to $D_0,D_2$ gives exactly such a coincidence.
\end{proof}

\begin{corollary}[\texttt{conjecture\_00000007608\_false}]
(U) is false. It already fails for the rotation by $90^\circ$, $\rho=\left(\begin{smallmatrix}0&-1\\
1&0\end{smallmatrix}\right)$.
\end{corollary}

\paragraph{Explicit integer examples.} Let $R_1=\operatorname{diag}(1,-1)$,
$R_2=\left(\begin{smallmatrix}0&1\\ 1&0\end{smallmatrix}\right)$,
$R_3=\operatorname{diag}(-1,1)$ and $R_4=\left(\begin{smallmatrix}0&-1\\
-1&0\end{smallmatrix}\right)$. Their mirrors are the $x$-axis, the line $y=x$, the $y$-axis and
the line $y=-x$; these are pairwise different (\texttt{four\_mirrors\_distinct}, witnessed by the
vectors $(1,0)$, $(1,1)$ and $(0,1)$). Then:
\begin{itemize}
\item $R_2R_1=R_4R_3=\rho_{90^\circ}$ (\texttt{rot90\_two\_decompositions});
\item $R_3R_1=R_4R_2=R_2R_4=-I=\rho_{180^\circ}$ (\texttt{rot180\_two\_decompositions});
\item $\rho_{90^\circ}A_2=\frac15\left(\begin{smallmatrix}-4&-3\\ -3&4\end{smallmatrix}\right)$
(\texttt{Aref2\_example}).
\end{itemize}

\section{Readings}\label{sec:readings}

\begin{itemize}
\item \textbf{Ordered vs.\ unordered.} Lean refutes the weakest literal form, \texttt{UniqueClause}:
two non-parallel decompositions should have the same \emph{unordered} pair of mirrors. The
ordered, matrix-level form \texttt{UniqueClauseOrdered} implies it
(\texttt{ordered\_implies\_unordered}), so it is false too (\texttt{conjecture\_false\_ordered}).
\item \textbf{Which rotation.} Non-uniqueness is not a special phenomenon of one rotation. It
holds for every rotation $\neq I$: Theorem~\ref{thm:main} over $\R$, and
\texttt{every\_nontrivial\_rotation\_fails} for all rational rotations in Lean. For $\rho=I$
every decomposition is $(A,A)$, with equal (``parallel'') mirrors, so (U) holds vacuously there
(\texttt{uniqueFor\_id}). This also shows that the formal predicate is not trivially false.
\item \textbf{Real vs.\ rational.} Lean works over $\Q$. Theorem~\ref{thm:main} is the real
statement, and \texttt{verify.py} illustrates it numerically for irrational angles.
\item \textbf{Affine reflections.} An affine reflection is $v\mapsto Mv+u$, with $M$ a linear
reflection, that is an involution; its mirror is an affine line. A rotation about a point $c$
is $\rho_c(v)=\rho(v-c)+c$. If two mirrors meet in a point $P$, the composition fixes $P$. If
they are parallel and distinct, the composition is a nonzero translation. For example, the
mirrors $x=0$ and $x=1$ give $v\mapsto v+(2,0)$, which has no fixed point
(\texttt{parallel\_gives\_translation}). So for $\rho_c\neq\id$ both mirrors pass through $c$,
and conjugating by the translation $v\mapsto v-c$ identifies the affine decompositions of
$\rho_c$ with the linear decompositions of $\rho$. The affine reading therefore has the same
infinite family. Concretely, the rotation by $90^\circ$ about $(1,1)$ equals
$S_{y=x}\circ S_{y=1}=S_{x+y=2}\circ S_{x=1}$, with four different non-parallel mirrors through
$(1,1)$ (\texttt{affine\_two\_decompositions}).
\item \textbf{$\R^n$ and hyperplane mirrors.} The definition speaks of hyperplane mirrors. In
$\R^n$ a product of two hyperplane reflections with distinct mirrors $H_1,H_2$ is a rotation in
the plane $(H_1\cap H_2)^\perp$ about the axis $H_1\cap H_2$. A rotation $\rho\neq I$ of this
kind (angle $\theta$ in a plane $\Pi$, fixing $\Pi^\perp$) is decomposed exactly by the pairs of
hyperplanes containing $\Pi^\perp$, i.e.\ $H_i=\ell_i\oplus\Pi^\perp$, where $(\ell_1,\ell_2)$
runs over the decompositions of $\rho|_\Pi$ in Theorem~\ref{thm:main}. Again there is a circle's
worth. In Lean, \texttt{r3\_two\_decompositions} embeds the $90^\circ$ example block-diagonally in
$\R^3$. There the mirrors are the planes $y=0$, $x=y$, $x=0$ and $x=-y$, each containing the
$z$-axis, and they are pairwise different. Parallel hyperplanes through the origin are equal,
so the exception again never applies.
\item \textbf{Point reflections.} If ``reflection'' meant a point reflection $v\mapsto2c-v$, the
composition of two of them would be a translation, never a nontrivial rotation, and (C1) would
fail outright. The definition (``hyperplane mirror structure'') excludes this reading.
\item \textbf{The charitable escape: ``unique up to rotating both mirrors''.} The classical
textbook statement is that the decomposition is unique \emph{up to rotating both mirrors
simultaneously about the centre}: only the angle between the mirrors, $\theta/2\bmod\pi$, is
determined. This is true, and it follows from Theorem~\ref{thm:main}(1): any two decompositions
$(R_\alpha,R_{\alpha+\theta/2})$ and $(R_{\alpha'},R_{\alpha'+\theta/2})$ differ by rotating both
mirrors by $\alpha'-\alpha$. We do \emph{not} refute this version. However, it is not what the
conjecture states. The conjecture asserts uniqueness and names parallel mirrors as \emph{the}
exception, and the Chinese text says plainly that the decomposition is unique. Under the
``up to rotation'' reading there would be no exception at all for $\rho\neq I$, and parallel
mirrors would not be an exception to uniqueness but a different case. The literal clause
therefore claims uniqueness of the mirror pair, and that is what we refute.
\item \textbf{Clause (C1).} For linear reflections (C1) is true: $BA$ is orthogonal with
$\det=(-1)^2=1$ (\texttt{refl\_mul\_refl}). With equal mirrors the composition is the
identity, the rotation by $0$. For affine reflections with parallel distinct mirrors, (C1) is
false: the composition is a translation with no fixed point. This is the exception that (C4)
mentions, so (C1) and (C4) together are correct. The disproof does not use them.
\end{itemize}

\section{Machine verification}

\paragraph{Lean 4 (\texttt{lean4/Main.lean}, core only, v4.19.0).} From scratch it defines:
\begin{itemize}
\item integer $2\times2$ matrices \texttt{M2} (product, transpose, determinant, action on
$\Z^2$), and rational matrices \texttt{QM} $=m/q$ with equality \texttt{QEq}, proved to be an
equivalence relation;
\item \texttt{IsRefl} (orthogonal, $\det=-1$, involution), \texttt{IsRot}, the fixed-vector
predicate \texttt{Fix}, \texttt{SameMirror} (parallel $=$ equal mirrors), and \texttt{Decomp};
\item the claims \texttt{UniqueFor}~$\rho$, \texttt{UniqueClause} and
\texttt{UniqueClauseOrdered}.
\end{itemize}
The main theorems are:
\begin{itemize}
\item \texttt{conjecture\_00000007608\_false : $\neg$ UniqueClause};
\item \texttt{conjecture\_false\_rot90\_explicit} (the same, directly from $R_2R_1=R_4R_3$);
\item \texttt{conjecture\_false\_ordered};
\item \texttt{every\_nontrivial\_rotation\_fails};
\item \texttt{infinitely\_many\_decompositions};
\item \texttt{decomp\_classification}, \texttt{refl\_mul\_refl}, \texttt{isRefl\_iff} and
\texttt{uniqueFor\_id};
\item the concrete $90^\circ$, $180^\circ$, $\R^3$ and affine examples.
\end{itemize}
A small tactic \texttt{ring\_z} (expand, AC-normalise, \texttt{omega}) proves the polynomial
identities, since core Lean has no \texttt{ring}. The concrete checks use \texttt{decide}. There
is no \texttt{sorry} and no \texttt{native\_decide}. The axioms used are at most
\texttt{propext} and \texttt{Quot.sound}.

\emph{Limitations.} Lean works over $\Q$, not $\R$; the real statement
(Theorem~\ref{thm:main}) is proved in this report. Mirrors are compared through integer
vectors, which loses nothing, since a rational vector is fixed iff an integer multiple of it is.
The $\R^3$ and affine parts are concrete integer examples (\texttt{decide}). The general $\R^n$
and affine reductions above, and the ``up to rotation'' statement, are proved only in this
report.

\paragraph{Python (\texttt{verify.py}, standard library).} A different method: exact
\texttt{Fraction} matrices, mirrors computed as kernels by Gaussian elimination, a small
polynomial engine for the symbolic identities, and brute force. For 123 nontrivial rational
rotations and 86 rational parameters $t$, every $(A_t,\rho A_t)$ is a non-parallel decomposition
with pairwise different mirrors. In a set $S$ of 124 rational reflections it enumerates all
pairs $(A,B)\in S^2$ with $BA=\rho$. For $\rho_{90^\circ}$ there are 124 such pairs, all
non-parallel, with 124 different unordered mirror pairs. For $\rho=I$ all of them are $(A,A)$.
It also checks the $\R^3$ and affine examples, and evaluates
$R_{\alpha+\theta/2}R_\alpha=\rho_\theta$ in floating point for irrational angles.

\end{document}
